%! Characteristics of Exponential Functions — Unit 6, Lesson 6.0 — Warm-Up
\documentclass[10pt]{article}
\usepackage{algebra2-article}
\usepackage{algebra2-boxes}
\newcommand{\TallMath}[1]{$\displaystyle #1\rule[-1.4em]{0pt}{3.2em}$}

\begin{document}

\pageheader{Unit 6, Lesson 6.0}{Warm-Up}
\namedateperiod

\vspace{0.10in}

\begin{notesbox}{Getting Ready: Skills We'll Use Today}
\small These three skills power today's lesson on reading exponential graphs. Work quickly.

\vspace{4pt}
\textbf{1. Powers, including the strange exponents.} Evaluate each one. Give exact fractions.
\begin{center}
\renewcommand{\arraystretch}{1.5}
\begin{tabular}{@{}l l l l@{}}
(a) \; $2^{3}=$ \blank{1.6cm} & (b) \; $2^{1}=$ \blank{1.6cm} &
(c) \; $2^{0}=$ \blank{1.6cm} & (d) \; $2^{-1}=$ \blank{1.6cm} \\
(e) \; $2^{-2}=$ \blank{1.6cm} & (f) \; $2^{-4}=$ \blank{1.6cm} &
(g) \; $2^{-8}=$ \blank{1.6cm} & (h) \; $5^{0}=$ \blank{1.6cm} \\
\end{tabular}
\end{center}
Look at (d) through (g). As the exponent gets more and more negative, the answers get
\blank{2.6cm}. Could $2^{x}$ ever be \emph{exactly} $0$? Explain. \writeline

\vspace{4pt}
\textbf{2. Two tables, two different patterns.} Fill in the last entry of each, then say what you
\emph{do} to one $y$-value to get the next one.
\begin{center}
\renewcommand{\arraystretch}{1.3}
\begin{tabular}{|c|c|c|c|c|c|}
\hline
\multicolumn{6}{|c|}{\small Table A} \\ \hline
$x$ & $0$ & $1$ & $2$ & $3$ & $4$ \\ \hline
$y$ & $2$ & $5$ & $8$ & $11$ & \blank{0.8cm} \\ \hline
\end{tabular}
\hspace{0.2in}
\begin{tabular}{|c|c|c|c|c|c|}
\hline
\multicolumn{6}{|c|}{\small Table B} \\ \hline
$x$ & $0$ & $1$ & $2$ & $3$ & $4$ \\ \hline
$y$ & $2$ & $6$ & $18$ & $54$ & \blank{0.8cm} \\ \hline
\end{tabular}
\end{center}
Table A: each $y$ is the one before it \blank{4.0cm}. \quad
Table B: each $y$ is the one before it \blank{4.0cm}.

\vspace{4pt}
\textbf{3. A rumor.} Three people hear a rumor at hour $0$. Every hour after that, the number of
people who have heard it \textbf{doubles}. Complete the table.
\begin{center}
\renewcommand{\arraystretch}{1.3}
\begin{tabular}{|c|c|c|c|c|c|}
\hline
Hour $t$ & $0$ & $1$ & $2$ & $3$ & $4$ \\ \hline
People & $3$ & \blank{1.0cm} & \blank{1.0cm} & \blank{1.0cm} & \blank{1.0cm} \\ \hline
\end{tabular}
\end{center}
How many people had heard it at hour $0$? \blank{1.6cm} \quad What are you multiplying by each
hour? \blank{1.6cm} \\
Which of the two tables in item 2 does this rumor behave like? \blank{2.0cm}

\end{notesbox}

\end{document}
